\documentclass[]{../../../../tex/classes/styledArticle}
\usepackage{../../../../tex/packages/hebrewSupport}
\usepackage{../../../../tex/packages/mathShortcuts}
\usepackage{../../../../tex/packages/theoremsSupport}

\author{שחר פרץ}
\title{חדו''א 2א $\sim$ \textit{הרצאה 19}}
\begin{document}
	\maketitle
	\textbf{מרצה: }יעל קרשון
	
	בשיעור שעבר הוכחנו את משפט פייר (Fejer) – עבור $f \in C(\tb)$ מתקיים ש־$\sg_N f \tou f$ כאשר $\sg_N = \frac{1}{N + 1} \sum_{n = 0}^{N}S_nf$. קיבלנו כמה מסקנות מהמשפט: 
	\begin{enumerate}[(A)]
		\item כל $f \in C(\tb)$ ניתנות לקירוב במ''ש ע''י פולינומים טריגונומטריים, כלומר לכל $\eg > 0$ קיים פולינום טריגונומטרי $p$ כך ש־$\max\sof{f - p} < \eg$. 
		\item השתמשנו ב(א) כדי להוכיח התכנסות ב־$L_2$, כלומר לכל $g \in R(\tb)$ טור פורייה של $g$ שואף ל־$g$ ב־$L_2$. 
		\item עבור $f, g \in C(\tb)$, אם $\forall n \in \N \co \hat f(n) = \hat g(n)$ אז $f = g$. 
	\end{enumerate}
	\begin{proof}[הוכחת (ג) תוך שימוש במשפט פייר]
		לכל $x$, מפייר (הפונקציות רציפות) $(\sg_N f)(x) \toinf f(x)$ וכן $(\sg_N g)(x) \toinf g(x)$ נקודתית, כי התכנסות במ''ש גוררת התכנסות נקודתית. נבחין ש־$(\sg_Nf)(x) = (\sg_Ng)(x)$, ולכן מיחידות הגבול של סדרת מספרים, נסיק ש־$f(x) = g(x)$. 
	\end{proof}
	
	\begin{proof}[הוכחה אחרת ל־(ג), תוך שימוש ביחידות הגבול ב־$L_2$]
		ידוע שבמובן $L_2$ מתקבל ש־$S_Nf \to f$ ו־$S_Ng \to g$, ואז מיחידות הגבול ב־$L_2$ בהכרח $\norm{f - g} = 0$. אך $f, g$. אך $f, g$ רציפות כלומר $f - g = 0$ ואז $f = g$ כדרוש. 
	\end{proof}
	
	\rmark{גם ההוכחה תוך שימוש ביחידות הגבול ב־$L_2$ מסתמכת על פייר, כי הוכחנו את התכנסות הגבול ב־$L_2$ באמצעות פייר}
	
	
	\textbf{מעתה ואילך מתחיל הפרק של מרחבים מטריים. }למדתי אותו כבר עצמאית, אז לא תראו סיכומים שלו יותר. אני כרגע עובד על סיכום לקורס טופולוגיה אז אולי בהמשך אכתוב עצמאית סיכום משלי לפרק הזה, שלא נסמך על הרצאות מהכיתה. 
	
	\ndoc
\end{document}